% ARCHIVED FROM v82 EXACT SOURCE EXCERPT: 02_fixed-h-mcfg/main.tex
% Source reference: archive/mcfg-v82-pretrim-full-20261018 (commit bcc0caf0eade66f552127171c0de24d2b92950b0)
% Restoration: replace the corresponding trimmed location with this excerpt (without these header lines).

\section{Relation to Yoshinaka's General Framework}
\label{sec:general-framework}

The bounded-rank learner is a typed instance of
Yoshinaka's substitutable-signature framework
\cite{Yoshinaka2015General}. We make the
correspondence explicit to distinguish prior
qualitative learning results from the total-length
bounds of Section~\ref{sec:char-data}.

Fix finite $\Sigma$, $f,r\ge1$ and
$h:\Sigma^*\to M$, and put $M^+:=h(\Sigma^+)$.
The signature $\mathfrak S_{f,r,h}$ has a
start sort $O_*=\Sigma^*$ and tagged sorts
\[
 T_{d,\boldsymbol\mu}
 =\{\tuple w\in(\Sigma^+)^d:
 h^{(d)}(\tuple w)=\boldsymbol\mu\},
 \qquad
 1\le d\le f,\quad\boldsymbol\mu\in(M^+)^d.
\]
Its function names comprise all well-typed
good MCFG templates of child rank at most $r$
(including rank zero); promotions
$T_{1,\mu}\to O_*$; and a nullary $\eps$
constant of sort $O_*$. No function takes an
$O_*$ input. After forgetting the type tags, this is a
finite sort refinement, in the sense of
\cite[Section~8]{Yoshinaka2015General}, of
the signature of good $\lambda$-free MCFG
operations on nonempty tuples. That
untyped signature is a nonempty-component
variant of \cite[Example~3]{Yoshinaka2015General},
rather than literally the signature of that
example; at fan-out one the prefix--suffix
case corresponds to
\cite[Example~4]{Yoshinaka2015General}.
Write $\mathcal S_O,\mathcal C_O$ for the
substructures and contexts realizable in
some grammar over this signature.

\begin{lemma}[Signature contexts are admissible sentence contexts]
\label{lem:signature-contexts}
For every non-start sort,
\[
 \mathcal S_{T_{d,\boldsymbol\mu}}
 =T_{d,\boldsymbol\mu},
 \qquad
 \mathcal C_{T_{d,\boldsymbol\mu}}
 =\{\tuple x\mapsto E[\tuple x]:E\in\mathcal E_d\}.
\]
The latter equality is between \emph{functions}
on $T_{d,\boldsymbol\mu}$.
Also $\mathcal C_{O_*}$ contains only the identity.
\end{lemma}

\begin{proof}
Every element of a non-start sort is a rank-zero
good template constant. No operation takes
$O_*$ as input. For a context of non-start
sort, induct along the path from its hole
to the start. Composing a good template with
an admissible parent context preserves the
order and single occurrence of the
distinguished child's variables.
Consecutive holes in one output component
are separated by a terminal or a nonempty
sibling component (nonmerging); holes in
different output components are separated
by a nonempty parent-context separator.
Thus every resulting sentence context
belongs to $\mathcal E_d$, by repeated
application of Lemma~\ref{lem:filling-identity}.

Conversely, for any
$E=v_0\blank_1v_1\cdots\blank_dv_d
\in\mathcal E_d$, the unary good template
$(v_0\xi_1^1v_1\cdots\xi_1^dv_d)$
with a rank-zero child and start promotion
realizes $\tuple x\mapsto E[\tuple x]$.
Here the internal $v_i$ are nonempty,
as required by goodness; the output type
is fixed by $h$.
\end{proof}

\begin{proposition}[Bounded-rank coverage by Yoshinaka's 2015 framework]
\label{prop:yoshinaka2015-embedding}
The signature $\mathfrak S_{f,r,h}$ satisfies
Assumptions~1--3 of
\cite{Yoshinaka2015General}, and
\(\mathsf{Subst}(\mathfrak S_{f,r,h}) =\Cmcf{f}{h}\cap\mathsf{MCFL}(f,r).\)
Consequently Yoshinaka's identification theorem
applies to the bounded-rank class
\cite[Theorem~1]{Yoshinaka2015General}:
for fixed $f,r,h$, updates are polynomial
in the accumulated sample, and characteristic
sets have cardinality at most
$|N^\sharp|+|P^\sharp|$ for a trimmed
signature presentation
\cite[Lemma~3]{Yoshinaka2015General}.
This does not bound the total \emph{length}
of those samples.
\end{proposition}

\begin{proof}
Every signature grammar is a fan-out-$\le f$,
rank-$\le r$ MCFG after forgetting tags.
Conversely, rank-preserving good normalization
(Proposition~\ref{prop:good-normalization}),
trimmed yield-type refinement
(Definition~\ref{def:output-refinement}),
start promotions and the separate $\eps$
constant present every language in
$\mathsf{MCFL}(f,r)$ over this signature.
At the start sort substitutability is
vacuous; at all other sorts
Lemma~\ref{lem:signature-contexts} identifies
the four-instance signature implication
with Definition~\ref{def:tuple-substitutable}
for each arity and componentwise type.
The trimmed presentation has no useless
nonterminals or functions, as required in
\cite[Section~6]{Yoshinaka2015General}.
This proves the equality.

For the effective assumptions, sort/context
membership and template evaluation are
polynomial-time by the lemma and finite
typing; uniform bounded-rank MCFG membership
is polynomial
\cite[Proposition~4]{Yoshinaka2011}.
All operations have arity at most $r$.
For a sentence of length $n$, enumerate
the start word, each tagged tuple of up
to $f$ nonoverlapping nonempty intervals
with nonempty internal gaps, and the
associated context functions. For each
parent tuple occurrence, enumerate the
rank-$\le r$ good templates determined by
ordered child-component intervals (at most
$fr$ such intervals) and the remaining
terminal material, recording the actual
\emph{function names}. Add promotions and
the $\eps$ constant as appropriate.
By Lemma~\ref{lem:composition-witness-materialization},
every realizable function instance occurs.
At most $2f(r+1)$ cut positions are needed,
so the enumeration is polynomial in $n$
for fixed $f,r,h$. These are precisely
Assumptions~1--3; applying Yoshinaka's
Theorem~1 and Lemma~3 gives the stated
learning consequences.
\end{proof}

\begin{lemma}[Substructure-indexed and guarded reconstruction]
\label{lem:substp}
For the preceding signature and enumerations,
Yoshinaka's $\mathrm{SUBST}_{\mathrm P}(K)$
and $\mathcal B^{\le r}_{f,h}(K)$ generate
the same language for every finite
positive sample $K\subseteq\Sigma^*$.
\end{lemma}

\begin{proof}
Identify their non-start symbols indexed
by the same observed tuples and typed sorts.
A $\mathrm{SUBST}_{\mathrm P}$ rule
with parent $s_0$ and template value
$z=\rho(s_1,\ldots,s_t)$ requires a
common $c\in\mathrm{CON}_O(K)$ such that
$c\odot s_0,c\odot z\in K$,
in the notation of
\cite{Yoshinaka2015General}. By Lemma~\ref{lem:signature-contexts}
this is an admissible sentence context.
The guarded Rule~(Link) sends $[s_0]$
to $[z]$, followed by Rule~(Comp), or
Rule~(Const) for $t=0$, whose witness
is materialized by
Lemma~\ref{lem:composition-witness-materialization}.
Conversely, every guarded Link is a
$\mathrm{SUBST}_{\mathrm P}$ instance of
the unary identity template; every Comp
and Const is its own witnessed parent
instance. Induction on derivations
therefore equates the non-start
symbol languages.

At the start sort the only context is
the identity and no operation accepts
a start-sort argument. Each observed
start symbol consequently has the
same language
\(\bigcup_{w\in K\cap\Sigma^+}L([(w)]) \ \cup\ (\{\eps\}\cap K),\)
by the promotions and direct $\eps$
constant. This is precisely guarded
Rule~(Start); for $K=\emptyset$ both
languages are empty.
\end{proof}

\paragraph{Characteristic witness correspondence.}
Yoshinaka's proof uses the context-indexed
$\mathrm{SUBST}_{\mathrm D}$, whereas its
Section~7 discusses $\mathrm{SUBST}_{\mathrm P}$
under the stated additional assumption
\cite{Yoshinaka2015General}.
The sample-level equality here is established
directly by Lemma~\ref{lem:substp}.
Choosing its non-start witnesses as
$c_X=\chi(X)$ and $s_X=\omega(X)$
gives precisely $\Wit(\widetilde G_0)$
at non-start sorts; promoted start
witnesses are among these chosen words,
and $\eps$ contributes its direct
datum when present. This identifies
the positive-word witness \emph{pattern}
with consistent start conventions,
not the syntactic hypothesis grammars.

